\documentclass[a4paper,12pt]{article}

\usepackage[utf8]{inputenc}
\usepackage[T1]{fontenc}
\usepackage{titlesec}
\usepackage{changepage} 
\usepackage[italian]{babel}
\usepackage{amsmath,amssymb,amsthm}
\usepackage{graphicx}
\usepackage[margin=2.5cm]{geometry}
\usepackage[hidelinks]{hyperref}
\usepackage{changepage}
\usepackage{tikz}
\usepackage{algorithm2e}
\usepackage{algpseudocode}
\usepackage{listings}

\usetikzlibrary{positioning}

\newenvironment{paragrafo}[1]{
    \par
    \noindent\textbf{#1}\par
    \begin{adjustwidth}{2em}{0pt}
    \setlength{\parindent}{0pt}
    \noindent\ignorespaces
}{
    \end{adjustwidth}
    \par
}

\newcommand{\memoria}[1]{%
\begin{tikzpicture}[
    cella/.style={
        draw,
        minimum width=0.8cm,
        minimum height=0.55cm,
        inner sep=0pt
    }
]
\foreach \val [count=\i from 0] in {#1}{
    \node[cella] (c\i) at (\i*1.05,0) {\val};
    \node[above=2pt of c\i] {\scriptsize \i};
}
\end{tikzpicture}
}

\newcommand{\memoriaunita}[1]{%
\begin{tikzpicture}[
    cella/.style={
        draw,
        minimum width=0.8cm,
        minimum height=0.55cm,
        inner sep=0pt,
        outer sep=0pt
    }
]
\foreach \val [count=\i from 0] in {#1}{
    \node[cella, anchor=west] (c\i) at ({0.8*\i},0) {\val};
    \node[above=2pt of c\i] {\scriptsize \i};
}
\end{tikzpicture}%
}

\title{\Huge Appunti di Algoritmi e Complessità\\\small per il corso di Boldi P.  @ Università degli studi di Milano}
\author{Manuele Ferro}  
\date{2026 / 2027}


\begin{document}

% Source - https://tex.stackexchange.com/a/741116
% Posted by satk0
% Retrieved 2026-09-23, License - CC BY-SA 4.0

\lstdefinelanguage{Pseudocode}{
  keywords={input, output, return, datatype, function, in, if, else, foreach, while, begin, end
  },
  sensitive=false,
  comment=[l]{//},
  morecomment=[s]{/*}{*/},
  morestring=[b]',
  morestring=[b]"
}

\definecolor{codegreen}{rgb}{0,0.6,0}
\definecolor{codegray}{rgb}{0.5,0.5,0.5}
\definecolor{codepurple}{rgb}{0.58,0,0.82}
\definecolor{backcolour}{rgb}{0.95,0.95,0.92}

\lstdefinestyle{mystyle}{
    backgroundcolor=\color{backcolour},   
    commentstyle=\color{codegreen},
    keywordstyle=\color{magenta},
    numberstyle=\tiny\color{codegray},
    stringstyle=\color{codepurple},
    basicstyle=\ttfamily\footnotesize,
    breakatwhitespace=false,         
    breaklines=true,                 
    captionpos=b,                    
    keepspaces=true,                 
    numbers=left,                    
    numbersep=5pt,                  
    showspaces=false,                
    showstringspaces=false,
    showtabs=true,                  
    tabsize=2
}


\maketitle
\tableofcontents
\newpage

\section{Introduzione}
\section{Analisi di algoritmi}
\subsection{Load balancing}
\begin{paragrafo}{Problema:}
    \begin{paragrafo}{Input:}
        Ho in input un numero $n \in \mathbb{N}^+$ di task e un numero $m \in \mathbb{N}^+$ di macchinari.

        Ognuna di queste task ha un tempo di esecuzione $t_0,t_1,\dots t_{n-1} \in \mathbb{N}^+$

        Indichiamo con $L_i$ il carico associato a una macchina $i$.
        
        Indichiamo con $L$ il costo, definito come  $L = \max_i{L_i}$
    \end{paragrafo}

    \begin{paragrafo}{Soluzioni ammesse:}
        Le soluzioni ammesse per questo problema sono delle partizioni di $n$ per $m$:

        \begin{center}
            $A_0,A_1,\dots ,A_{m-1}$
        \end{center}

        Il problema è di ricerca del minimo.
    \end{paragrafo}

    \begin{paragrafo}{Esempio grafico:}
        Supponiamo di avere 8 task, e 3 macchinari su cui eseguirle.
        
        \textbf{Tasks:} 
        \begin{center}
            \memoria{3,2,3,1,1,4,5,1}
        \end{center}
        \smallskip
        Possiamo trovare una possibile partizione delle task per i macchinari del tipo:

        \smallskip
        \begin{center}
            \begin{tabular}{r@{\qquad}l@{\qquad}r}
                $m$ & tasks & $L_i$ \\
                0 & \memoriaunita{3,1,1,2} & 7 \\
                1 & \memoriaunita{5,1,3}   & 8 \\
                2 & \memoriaunita{4,1}     & 5
            \end{tabular}
        \end{center}
        \smallskip
        Vediamo facilmente che il costo $L$ è il massimo tra i 3, ovvero $8$.
        Possiamo allo stesso tempo affermare che non è la soluzione ottima, infatti, è semplice evidenziarne una migliore.
        
        \smallskip
        \begin{center}
            \begin{tabular}{r@{\qquad}l@{\qquad}r}
                $m$ & tasks & $L_i$ \\
                0 & \memoriaunita{5,1,1} & 7 \\
                1 & \memoriaunita{4,3}   & 7\\
                2 & \memoriaunita{3,2,1}     & 6
            \end{tabular}
        \end{center}
        \smallskip

    \end{paragrafo}
    
    \break

    \begin{paragrafo}{Teorema.} 
        Il problema Load Balancing è NPO Completo

        \hspace*{2em} Questo ci porta subito a studiare le soluzioni approssimate.
    \end{paragrafo}
        
    Inserire algoritmo qui che non ho sbatti, e anche i vari tempi e cazzate varie.

    \begin{paragrafo}{Teorema.} 
        Greedy Load Balancing è 2 approssimato per Load Balancing.
    \end{paragrafo}
    \begin{paragrafo}{Dimostrazione}
        \textbf{Osservazione 1.} $L^* \ge \frac{1}{m}\sum_jt_j$
        
        \textbf{Dim 1.} $\sum_{i} L_i^* = \sum_{j} t_j$
        
        \hspace*{2em} Per pidgeonholing, non è possibile che $\forall i L^*_i < \frac 1 m \sum_{j}t_j$

        \hspace*{2em} quindi

        \hspace*{2em} $\exists i L^*_i \ge \frac 1 m \sum_j  t_j$

        \textbf{Osservazione 2.} $L^* \ge \max_j t_j$

        \vspace*{2em}

        Sia $\hat i$ la macchina più carica alla fine dell'esecuzione e $\hat j$ l'ultima task assegnata a $\hat i$.

        Allora possiamo dire che:

            $$L_{\hat i} - t_{\hat j} \le L_i'\ \forall i \in m \le L_{i}$$

        A questo punto sommiamo le disequazioni:

            $$L_{\hat i} - t_{\hat j} \le L_i'$$
            $$m(L_{\hat i} - t_{\hat j}) \le \sum_i L_i$$
            $$L_{\hat i} - t_{\hat j} \le \frac 1 m \sum_i L_i = \frac 1 m \sum_j t_j\le L^*$$
            $$L = L_{\hat i} = (L_{\hat i} - t_{\hat j}) + t_{\hat j} \le L^* + t_{\hat j}\le L^*+L^*$$
            $$L\le2L^*$$

        \textbf{Corollario.} Load Balancing $\in$ 2-APX
            
            

    \end{paragrafo}

    \begin{paragrafo}{Teorema.} 
        $\forall \epsilon > 0$ esiste un input in cui Greedy Load Balancing produce una soluzione di costo:
        $$2-\epsilon \le \frac{L}{L^*} \le 2$$
    \end{paragrafo}
    \begin{paragrafo}{Dimostrazione}
        Prendiamo $m > \frac 1 \epsilon, \text{e } n = (m-1)+1$, con i primi $m-1$ input pari a $1$ el'ultimo pari a $m$.
        
        Applicando l'algoritmo su questi input otteniamo:
        \begin{center}
            \begin{tabular}{r@{\qquad}l@{\qquad}r}
                $m$ & tasks & $L_i$ \\
                0 & \memoriaunita{1,1,1} & 13 \\
                1 & \memoriaunita{1,1,3}   & 3 \\
                2 & \memoriaunita{1,1}     & 3
            \end{tabular}
        \end{center}
        $$L = 2m-1$$
        Notiamo invece il caso ottimale su questo input:
                \begin{center}
            \begin{tabular}{r@{\qquad}l@{\qquad}r}
                $m$ & tasks & $L_i$ \\
                0 & \memoriaunita{1,1,1} & 3 \\
                1 & \memoriaunita{1,1,1}   & 3 \\
                2 & \memoriaunita{1,3}     & 4
            \end{tabular}
        \end{center}
        $$L = 4, \text{ ottimo}$$ 
        \begin{paragrafo}{Piccola digressione}
            L'input definito è inoltre uno dei pochi casi in cui la media equivale al caso ottimo.
            $$\frac{m(m-1)+m}{m}= m-1+1 = \lceil m \rceil$$
        \end{paragrafo}

        $$\frac L {L^*} = \frac {2m-1} {m} = 2 - \frac 1 m \ge 2-\epsilon$$

    \end{paragrafo}

    Se ordino Greedy Load Balancing, perdo la sua proprietà di essere ONLINE, ma ne guadagno in approssimazione, infatti:

    \begin{lstlisting}[style=myStyle,mathescape=true,caption=Offline Greedy Load Balancing,label={Offline Greedy Load Balancing}]
    Sorto i task in ordine non crescente di durata
    Applico Greedy Load Balancing  
\end{lstlisting}

    \begin{paragrafo}{Teorema.}
        Offline Greedy Load Balancing è $\frac{3}{2}$ approssimato per Load Balancing.
    \end{paragrafo}

    \begin{paragrafo}{Dimostrazione.}
        \textbf{Osservazione 1.} Se $n \le m$, troviamo l'ottimo

        \hspace*{2em} Assumiamo quindi $n > m$

        \textbf{Osservazione 2.} $L^* \ge 2t_m$ dato che
            \begin{paragrafo}{Dim 2.}
                $\begin{underbrace} {t_0 \ge t_1 \ge \dots \ge t_{m-1}} \end{underbrace}_\text{m task} \ge t_m \dots \text{ sempre per pidgeonholing}$  
            
                $\exists i\ L^* \ge L^*_i \ge t_{1} + t_{2} \ge t_m + t_m$

                quindi
            
                $L^* \ge 2t_m$
            \end{paragrafo}
            
            Sia $\hat i$ la macchina più carica alla fine dell'esecuzione.

            Possiamo assumere che ad essa siano assegnati almeno due task, altrimenti avremmo l'ottimo.

            Sia $\hat j$ l'ultimo task assegnato:

            $$\hat j \ge m$$

            quindi

            $$t_{\hat j} \le t_m \le \frac 1 2 L^*$$

            quindi
            
            $$L = L_{\hat i} = L_{\hat i} - t_{\hat j} + t_{\hat j} \le \frac 3 2 L^*$$


            \textbf{N.B. 1:} L'analisi è migliorabile, dimostrando che l'algoritmo è $\frac 4 3 $ approssimato.

            \hspace*{2em} [Graham 1969] 

            \textbf{N.B. 2:} Load Balancing $\notin $ FPTAS (assunto $P \not = NP$)
            
            \hspace*{2em} [Hochbaum-Shmops 1988] 


    \end{paragrafo}
    


\end{paragrafo}

\end{document}
